% ===========================================================================
\section{Related Work}
\label{sec:related}
% ===========================================================================

\subsection{Global products and what they can and cannot resolve}

EDGAR \citep{edgar2025booklet,crippa2020} reports a Pakistani waste total covering disposal, burning
and wastewater together, and is the closest external comparator to the scope reported here.
Climate TRACE \citep{climatetrace} takes a different approach, inferring emissions from
asset-level and remote-sensing evidence. Both are global products: they apply a consistent
methodology across every country, which is what makes them comparable and also what limits
them nationally. A global waste model must assume a collection rate and a sewer condition for
Pakistan; a national compiler can consult the sanitation statistics and the district
population panel that describe it.

The disagreement between products is not a defect of any of them, but it is large enough to
matter. This inventory sits at \SI{141}{\percent} of EDGAR's year-matched Pakistani waste
figure, and
Section~\ref{sec:validation} argues that a single assumption---the sewer methane correction
factor---accounts for most of the difference.

\subsection{Waste generation and composition in South Asia}

Per-capita waste generation and its composition are the two activity inputs the sector cannot
avoid, and both are surveyed rather than metered. The global compilation by
\citet{kaza2018} is the standard reference; its Pakistan generation row re-bins the one
primary national survey, the Pak-EPA eight-city study \citep{pakepa1996}, whose rates are
used here directly, and its collection coverage supplies the uncollected fractions.
Composition matters more than it first appears: the fossil carbon fraction of municipal
waste---plastics, textiles, rubber---determines how much of the \COtwo{} from open burning is
reportable, because the biogenic remainder is a memo item under IPCC's convention rather than
a national-total emission.

\subsection{First Order Decay and its implementation choices}

Landfill methane is not a function of this year's waste but of every year's waste still
decaying, which is why IPCC's Volume~5 method is a first-order decay model over a deposit
history rather than an emission factor applied to an annual tonnage. Three implementation
choices then do most of the work and are reported inconsistently in the literature: the decay
constant, the methane correction factor describing site management, and how far back the
deposit series is extended before the model is considered spun up.

The third is the one this paper treats explicitly. A model started at the first year of an
activity panel implicitly asserts that nothing was landfilled before it, which is never true
and always biases the estimate low. The size of that bias is rarely reported.
Section~\ref{sec:results} reports it here.

\subsection{The wastewater methane correction factor}

For domestic wastewater the corresponding controversy is the methane correction factor, which
translates a sanitation pathway into the fraction of its organic load that degrades
anaerobically. IPCC's default table distinguishes stagnant sewers from flowing ones by an
order of magnitude in that fraction, and the distinction was drawn for systems that are
clearly one or the other. Open urban drains that are stagnant in the dry season and flowing
after rain, and which discharge to irrigation channels, are not obviously either.

Pakistani sanitation coverage itself is well characterised: the WHO/UNICEF Joint Monitoring
Programme \citep{jmp2023} reports the population shares using sewers, septic systems, latrines
and open defecation, and those shares are what allocate the national organic load across
pathways here.
What the JMP does not report, and no Pakistani source does, is the hydraulic condition of the
sewered fraction---which is the quantity the correction factor is asking about.
